%! From Combinations to Solutions --- Setting Up Ax = b — Unit 2, Lesson 2.0 — Lesson Plan
\documentclass[10pt]{article}
\usepackage{linalg-article}
\usepackage{linalg-boxes}
\usepackage{graphicx}   % -article does NOT load graphicx; needed for thumbnails/screenshots
\graphicspath{{images/}}
\newcommand{\TallMath}[1]{$\displaystyle #1\rule[-1.4em]{0pt}{3.2em}$}
\newcommand{\xx}{\mathbf{x}}
\newcommand{\bb}{\mathbf{b}}
\newcommand{\UnitNumberName}{Unit 2: Solving Linear Equations Ax = b \quad}
\newcommand{\LessonNumberName}{Lesson 2.0: From Combinations to Solutions --- Setting Up Ax = b}

\begin{document}
\begin{center}
    {\Large\bfseries \CourseName: \SchoolYear \\
    {\small\bfseries \UnitNumberName \LessonNumberName}}
\end{center}

\begin{tcolorbox}[colback=blush, colframe=burgundy, boxrule=0.9pt, arc=2mm,
    left=2mm, right=2mm, top=1.5mm, bottom=1.5mm]
    \textbf{Primary Objective:} Students can read the matrix equation $A\xx=\bb$ two ways ---
    as a \emph{combination of the columns} of $A$ and as a \emph{system of linear equations} ---
    solve a $2\times 2$ system by eliminating a variable, check the answer by rebuilding $\bb$,
    and explain what one, none, or infinitely many solutions means both geometrically (two lines)
    and in terms of whether $\bb$ is reachable.
\end{tcolorbox}

\renewcommand{\arraystretch}{1.4}
\begin{skillbox}[Priority Ideas \& Skills]{goldbox}
\begin{tabularx}{\linewidth}{@{}X|X@{}}
\textbf{Priority Ideas \& Skills} & \textbf{Key Understandings (the ``why'')} \\[2pt]
\hline
\vspace{-6pt}
\begin{itemize}[leftmargin=1.1em, nosep, after=\vspace{2pt}]
  \item See $A\xx=\bb$ three ways: matrix form, column combination, and a system of equations.
  \item Solve a $2\times 2$ system by subtracting to eliminate a variable, then back-substitute.
  \item Check a solution by substituting back and rebuilding $\bb$.
  \item Read solvability from the geometry (two lines) and from reachability.
\end{itemize}
&
\vspace{-6pt}
\begin{itemize}[leftmargin=1.1em, nosep, after=\vspace{2pt}]
  \item Unit~1 built $\bb$ \emph{from} $\xx$; Unit~2 runs it backwards --- find the $\xx$ that hits $\bb$.
  \item Eliminating a variable is the engine of the whole unit.
  \item The columns view and the rows view are one problem in two languages.
  \item One point / parallel / same line $=$ one / no / infinitely many solutions.
\end{itemize}
\end{tabularx}
\end{skillbox}

\begin{skillbox}[{Vocabulary, Concepts, \& Theorems}]{greenbox}
\renewcommand{\arraystretch}{1.3}
\begin{tabularx}{\linewidth}{@{}l X@{}}
\textbf{System of linear equations} & Several linear equations sharing the same unknowns; a solution must satisfy them all at once. \\
\textbf{Matrix form $A\xx=\bb$} & \TallMath{\begin{bsmallmatrix} 1 & 2 \\ 3 & 2 \end{bsmallmatrix}\begin{bsmallmatrix} x_1 \\ x_2 \end{bsmallmatrix}=\begin{bsmallmatrix} 5 \\ 11 \end{bsmallmatrix}}; the same system packed into one equation. \\
\textbf{Coefficient matrix $A$} & Its \emph{columns} are the base vectors being combined; its \emph{rows} give the equations. \\
\textbf{Solution} & The values of the unknowns $x_1, x_2$ that make every equation true. \\
\textbf{Elimination} & Subtract a multiple of one equation from another to remove a variable (the focus of 2.1). \\
\textbf{Consistent / inconsistent} & Has at least one solution / has no solution. \\
\textbf{Unique vs.\ infinitely many} & Lines cross once (one solution) vs.\ lie on top of each other (infinitely many). \\
\end{tabularx}
\end{skillbox}

\begin{fixedskillbox}[Activate Prior Knowledge \& Spiral Review]{blush}
    The Warm-Up rehearses the three skills this lesson stands on:
    \textbf{(1)} solving a small system by any Algebra-1 method (elimination will formalize it),
    \textbf{(2)} computing $A\xx$ as a \emph{combination of the columns} (Lesson~1.3), and
    \textbf{(3)} deciding whether a target $\bb$ is \emph{reachable} --- in the column space
    (Lesson~1.3). Debrief item~2 out loud: the product $\begin{bsmallmatrix} 5\\11 \end{bsmallmatrix}$
    is the very target students will solve \emph{back} to during the lesson.
\end{fixedskillbox}

\begin{skillbox}[Hook]{blush}
    Project two trail-mix ``bases.'' One scoop of \textbf{Base~A} is $1$ cup nuts and $3$ cups
    raisins; one scoop of \textbf{Base~B} is $2$ cups nuts and $2$ cups raisins. A recipe calls
    for exactly $5$ cups nuts and $11$ cups raisins. Ask: \textbf{``How many scoops of each base
    hit the target exactly --- and is there only one answer?''} Let a few students guess-and-check
    before naming the structure: this is $A\xx=\bb$ with the bases as \emph{columns}, run
    \emph{backwards} to find the scoops $\xx$.
\end{skillbox}

\begin{skillbox}[Lesson]{blush}
\begin{multicols}{2}
\textbf{1. Running Unit~1 backwards.} In Lesson~1.3 we \emph{built} $\bb=A\xx$: given a matrix
$A$ and weights $\xx$, combine the columns. Now flip it: given $A$ and a target $\bb$, find the
weights $\xx$ that \textbf{solve} $A\xx=\bb$. That is the question of all of Unit~2.

\textbf{2. Two languages, one problem.} Write the trail-mix problem both ways. \emph{Column
view:} $x_1\begin{bsmallmatrix} 1\\3 \end{bsmallmatrix}+x_2\begin{bsmallmatrix} 2\\2 \end{bsmallmatrix}=\begin{bsmallmatrix} 5\\11 \end{bsmallmatrix}$. \emph{Row view:} one equation per
row, $x_1+2x_2=5$ and $3x_1+2x_2=11$. Stress that these are the \emph{same} numbers, read
down the columns vs.\ across the rows.

\columnbreak

\textbf{3. Solve by elimination.} Subtract the first equation from the second to \textbf{kill}
$x_2$: $(3x_1+2x_2)-(x_1+2x_2)=11-5$ gives $2x_1=6$, so $x_1=3$; back-substitute to get $x_2=1$.
\textbf{Check} by rebuilding $\bb$: $3\begin{bsmallmatrix} 1\\3 \end{bsmallmatrix}+1\begin{bsmallmatrix} 2\\2 \end{bsmallmatrix}=\begin{bsmallmatrix} 5\\11 \end{bsmallmatrix}$. Name the move
\textbf{elimination} --- Lesson~2.1 makes it systematic.

\textbf{4. Geometry \& solvability.} Each row is a line; the solution is where they
\textbf{cross}. Lines crossing once $=$ one solution; \textbf{parallel} $=$ none (inconsistent);
the \textbf{same line} $=$ infinitely many. Tie back to 1.3: a solution exists exactly when
$\bb$ is \emph{reachable} (in the column space).

\textbf{5. The road ahead.} Preview Unit~2: \textbf{2.1} elimination in general, \textbf{2.2}
elimination matrices \& inverses, \textbf{2.3} $A=LU$, \textbf{2.4} permutations \& transposes.
\end{multicols}
\end{skillbox}

\begin{skillbox}[Active Monitoring]{redbox}
Circulate during the notes practice and Tier work. Watch for and correct:
\begin{itemize}[leftmargin=1.2em, nosep]
  \item mixing up the two readings --- combining down the columns vs.\ writing one equation per row;
  \item subtracting equations in the wrong order or forgetting to subtract the \emph{right-hand
        sides} too ($11-5$, not just the left sides);
  \item stopping at one variable and not back-substituting; skipping the check.
\end{itemize}
Cold-call prompts: ``What does $x_1$ mean in the trail-mix story?'' ``Which move made $x_2$
disappear?'' ``How would the picture look if there were \emph{no} solution?''
\end{skillbox}

\begin{skillbox}[Group Work \& Differentiation]{redbox}
\textbf{Blend the Batch} activity (tiered; mirrors the notes).
\begin{multicols}{3}
\textbf{Tier R --- Remediate.} One given system with an integer answer: rewrite $A\xx=\bb$ as
two equations, eliminate a variable, back-substitute, and check by rebuilding $\bb$.

\columnbreak
\textbf{Tier A --- Approaching.} A blend problem in context: set up $A\xx=\bb$ from two base
mixes and a target, solve it, and \emph{interpret} the scoops. Then a decide-the-outcome item.

\columnbreak
\textbf{Tier E --- Extension.} A system with \emph{no} solution and one with \emph{infinitely
many}; justify each from both the picture (two lines) and reachability of $\bb$.
\end{multicols}
\end{skillbox}

\begin{skillbox}[Individual Work \& Assessment]{redbox}
\textbf{Exit Ticket} (independent, no notes): rewrite a given $A\xx=\bb$ as a system of two
equations; solve it by eliminating a variable and check; and a \textbf{conceptual/justification
check} --- given two \emph{parallel} lines, explain why the system has no solution and connect it
to $\bb$ not being reachable. Collect and sort into ``got it / setup slip / elimination slip'' to
decide whether 2.1 opens with a quick re-teach.
\end{skillbox}

\begin{skillbox}[Reinforcement \& Extension]{goldbox}
\textbf{Homework:} translate between the three views of $A\xx=\bb$, solve $2\times 2$ systems by
elimination with a check, one applied blend problem with an interpret part, and a
one/none/infinitely-many classification. \textbf{Extension:} a $3\times 3$ system solved by
eliminating one variable at a time --- previews the staircase of 2.1. \textbf{Preview:}
Lesson~2.1, \emph{The Idea of Elimination} --- turning today's ``subtract to kill a variable''
into a reliable procedure with pivots.
\end{skillbox}
% ── Teacher notes (migrated from the component keys) ──
% One per component, titled for it. They live here rather than in the _key
% files so that every component is the same length blank and keyed.

\begin{teachernote}[Warm-Up]
Item~2 is deliberately the exact product $\begin{bsmallmatrix} 5\\11 \end{bsmallmatrix}$ students
will solve \emph{back} to in the lesson --- point that out in the debrief. Item~1 rehearses
elimination informally (adding the two equations cancels $y$); item~3 rehearses reachability from
Lesson~1.3, which becomes ``does $A\xx=\bb$ have a solution?''
\end{teachernote}

\begin{teachernote}[Guided Notes]
The whole lesson turns on one move: subtract to make a variable disappear. Keep the column view
and row view side by side so students see they are the same problem --- 2.1 formalizes the
subtraction into pivots. Common slip in \S3: subtracting the left sides but forgetting $11-5$ on
the right.
\end{teachernote}

\begin{teachernote}[Group Activity]
Tier A item~4 is the key interpret moment: a system can be \emph{mathematically} solvable yet
\emph{physically} impossible (negative scoops) --- this previews later work on when solutions are
usable. Tier E ties one/none/infinitely-many directly to reachability from Lesson~1.3: same
column line, but $\bb$ on it (infinitely many) vs.\ off it (none).
\end{teachernote}

\begin{teachernote}[Exit Ticket]
Item~3 checks the geometry--reachability link. Full credit needs both ``parallel lines don't
meet $\Rightarrow$ no solution'' \emph{and} the connection to $\bb$ not being in the column space.
\end{teachernote}

\end{document}
